\scanpage{18}
\mspage{15}

les deux éléments de $\delta I$, $I\in S$ (fixé) appartiennent à des classes distinctes \inserted{($\Longleftrightarrow$ idem pour les deux éléments des $B_x$, $x\in L$).}
\unsure{Je n'ai} pas la vérification ici. Si c'était faux, il faudrait le mettre en axiome. Sauf erreur, il doit être immédiat que \deleted{\unsure{il y a}} le nb de classes est $1$ ou $2$.

\deleted{On est orienté}

L'ens quotient $\underline{\omega}$ est l'ens des orientations de $\mathcal X$. \deleted{\unsure{On}} On a orienté $\mathcal X$ quand on a choisi $\omega\in\underline{\omega}$. Alors chaque $I$ a une origine, et pour chaque $x\in L$, on a un germe de segment \inserted{\uline{croissant}} issu de $x$ et d'origine $x$ (la \uline{branche sortante} de $x$).

Considérons alors la relation
\[
x<y\quad\overset{\mathrm{def}}{\Longleftrightarrow}\quad
\deleted{\unsure{x,y}}
\exists I\in S\text{ tel que }
\begin{array}{l}
x\in\operatorname{or}(I)\\
y\in\operatorname{ex}(I).
\end{array}
\]

Distinguons deux cas suivant que cette relation est transitive (\uline{cas linéaire}) ou non (cas circulaire).

I) Relation $x<y$ est transitive donc relation
\[
x\leq y\quad(\overset{\mathrm{def}}{\Longleftrightarrow}\ x=y\text{ ou }x<y)
\]
est un préordre,

\mspage{17}

\rotatebox{90}{\unsure{EXTRA STRONG}}
\rotatebox{90}{\unsure{EXTRA STRONG}}
